% Extracted from p06-testtime-scaling
\subsection{A Taxonomy of Test-Time Compute}
\label{ch-taxonomy}

\paragraph{Background}

\paragraph{Test-time scaling at large scale}

\citet{snell2024scaling} study how to spend a fixed inference FLOPs
budget across strategies: repeated sampling with majority vote or
best-of-$N$ selection, and sequential revision guided by verifiers. Two
findings matter here. First, for easier questions, parallel strategies
(best-of-$N$, self-consistency) are compute-optimal: independent samples
explore the answer space cheaply. Second, for harder questions,
sequential strategies win: iterative revision steered by a verifier
outperforms parallel sampling per FLOP. The crossover depends on task
difficulty and verifier quality. Both findings assume verifiers worth
having---a safe assumption at 8B, a loaded one at 44M.

\paragraph{Sampling-based methods}

Best-of-$N$ draws $N$ independent completions and keeps the best
according to a scorer. Self-consistency \citep{wang2023selfconsistency}
draws $N$ reasoning paths and takes the majority answer, needing no
learned scorer at all. Both multiply inference cost by $N$ (modulo
batching) and both live or die by selection: best-of-$N$ needs a
verifier better than chance; self-consistency needs the model's modal
answer to be the right one, which holds for reasoning tasks and fails
for open-ended generation.

\paragraph{Verifier-based methods}

Outcome verifiers score final answers; process verifiers score
intermediate steps \citep{cobbe2021verifiers}. They enable search
(beam search, MCTS, revision loops) and reranking. The large-scale
literature treats the verifier as a sunk cost---train once, reuse
forever. At small scale the verifier is a second model to train, store,
and serve, and its quality is bounded by the same data and capacity
constraints as the policy it judges.

\paragraph{Recurrent depth}

Recurrent transformers iterate a shared block for a fixed number of
loops during training \citep{dehghani2019universal}. Architectures in
this family, including OpenMythos \citep{gomez2024openmythos}, invite an
obvious test-time trick: run \emph{more} loops at inference than were
used in training. Unlike sampling, this costs no extra parameters and
no verifier. Whether the recurrent dynamics actually improve with extra
iterations is an empirical question, which Section~\ref{sec:loopsweep} answers for one
concrete model.

\paragraph{A three-axis taxonomy}

We classify test-time methods along three axes, each with a cost model
in units of base-model forward passes per token, $F$.

\paragraph{Axis I: sampling.} Best-of-$N$, self-consistency, majority
vote. Cost $\approx N \cdot F$, reducible by batching but not
asymptotically. Requires a selection mechanism: a verifier (Axis III) or
a majority signal. Failure mode: when the verifier is weak, best-of-$N$
selects confidently wrong answers---it amplifies verifier error, not
just policy error.

\paragraph{Axis II: depth.} Extra recurrent loops, adaptive computation
time, early-exit inversion (running \emph{past} the trained depth).
Cost $\approx L \cdot F$ for $L$ loops. Requires nothing but the model
itself. Failure mode: the dynamics may be calibrated to the training
iteration count, in which case extra loops add compute without adding
quality---or subtract it.

\paragraph{Axis III: verification.} Learned outcome or process reward
models, rerankers, revision critics. Cost $= V + N \cdot F$ where $V$ is
the verifier's forward cost per candidate. This axis rarely stands
alone; it is the steering wheel for Axes I and II. Failure mode: at
small scale, $V$ is not negligible, and a weak verifier mis-steers.

The axes compose: Snell's sequential revision is Axis II$\times$III;
best-of-$N$ with a reward model is Axis I$\times$III. The taxonomy's
point is to make the cost of each composition explicit before asking
which compositions a 44M budget can afford.

\paragraph{What changes below a billion parameters}

\paragraph{The verifier overhead problem}

Consider a learned verifier at half the policy's size---generous by
large-scale standards, where verifiers are often far smaller
relative to the policy. For an 8B policy, a 4B verifier is a 50\%
serving overhead. For a 44M policy, a 22M verifier is also a 50\%
overhead, but the absolute quality of a 22M verifier trained on the
same limited data is far lower: verifier quality scales with capacity
too, and there is no large-scale teacher to distil from cheaply. The
relative cost is the same; the absolute quality is not. This breaks
the large-scale assumption that verification is cheap and good.

\paragraph{Sampling arithmetic}

Best-of-8 on a 44M model spends $8 \times 44\text{M} = 352\text{M}$
parameters' worth of forward compute per token---roughly the cost of a
single greedy pass from a 350M model. The honest comparison is therefore
not ``44M vs.\ 44M with best-of-8'' but ``best-of-8 from 44M vs.\
greedy from 350M.'' If the 350M model wins that comparison, sampling is
a worse use of the FLOPs than training a bigger model. Whether it does
is task-dependent and, at small scale, largely unmeasured---which is
itself the point: the large-scale literature rarely runs this
comparison because at 70B the alternative (a 560B model) does not exist.

Self-consistency escapes the verifier problem but keeps the sampling
cost, and its majority-vote mechanism assumes the model's modal answer
is correct. On TinyStories-style generation, where there is no single
right answer, the mechanism has nothing to grip.

\paragraph{Depth: measured, not assumed}

Here we have direct evidence. Our 44M-parameter recurrent transformer
(43,974,722 parameters; dim 512; 6 prelude $+$ 2 coda layers; 4 training
loops; 8 experts; GQA; trained on TinyStories for 120M tokens, loss
9.5242 $\rightarrow$ 1.5263) was evaluated with loop counts swept from 1
to 16 over 256 held-out sequences (full results in
Section~\ref{sec:loopsweep}, Table~\ref{tab:sweep}).

Recurrence helps enormously up to the training depth---perplexity falls
by a factor of 6.7 from 1 to 4 loops---and then the gains stop dead.
Eight loops are slightly \emph{worse} than four; sixteen add nothing.
Generations remain coherent at all depths, so this is not collapse; it
is saturation. The recurrent dynamics are calibrated to the iteration
count they were trained with. Test-time depth extrapolation, the
cheapest axis in the taxonomy, is not free after all---at least not
without training-time support such as randomised loop counts.

\paragraph{The memory wall}

A final small-scale constraint: test-time methods that hold $N$
candidates or deep activation caches multiply memory, not just FLOPs.
At 44M the weights fit comfortably anywhere, which tempts wasteful
inference schemes; but the deployment targets for small models are
phones and browsers, where memory bandwidth---not capacity---is the
bottleneck, and $N$-way sampling spends it $N$ times over.

\paragraph{Position: three bets for small-scale test-time compute}

\begin{enumerate}
\item \textbf{Bet on trained-depth recurrence, not extrapolated depth.}
Running the loops you trained is pure win (our 1$\rightarrow$4 loop
result); running more is wishful thinking without training-time
support. Research should target \emph{adaptive} depth within the
trained range---early exit for easy tokens---rather than extrapolation
beyond it.

\item \textbf{Bet on cheap, honest verifiers.} Learned reward models are
too expensive and too weak at this scale. Heuristic verifiers---length
penalties, repetition detectors, grammar checkers, task-specific
validators---cost near zero and fail transparently. A verifier you can
read is better than a verifier you can only hope is calibrated.

\item \textbf{Sample only with a trustworthy selector.} Best-of-$N$ is
justified at small scale only when the selection signal is strong:
exact-match verifiers for constrained tasks, or self-consistency on
tasks with verifiable answers. Unverifiable open-ended generation
should not be sampled---the selector cannot tell good from lucky.
\end{enumerate}

Open questions: can randomised training loop counts unlock genuine
extrapolation? At what parameter count does a learned verifier become
worth its serving cost? And does the Snell easy/hard crossover replicate
when both policy and verifier are small?

\paragraph{Limitations}

This is a position paper; the taxonomy is argued, not experimentally
validated as a whole. The depth-axis evidence is a single model on a
single corpus---one careful data point, not a law. Snell's findings may
transfer better than we argue, particularly the easy-question sampling
results. We also ignore fine-tuning-time analogues (e.g., distilling
test-time compute back into weights), which may be the most practical
small-scale strategy of all.

\paragraph{Conclusion}

Test-time compute is real, but its economics are scale-dependent.
Sampling multiplies cost against a fixed budget, learned verifiers cost
a second model the budget cannot spare, and recurrent depth---the one
free axis---saturates at the training depth. The small-scale research
agenda should stop importing large-scale recipes and start designing
test-time methods for the regime where the verifier must cost less than
the model it serves.
